\cotefile{78}
\scanpage{32}
\deleted{et enfin la condition d'unimodularité}

\deleted{$4^\circ)$ $u_E$ unimodulaire}

% vertical rule beside the definition
On dit que la \deleted{\unsure{repr.}}\inserted{réalisation} régulière est \uline{unimodulaire} si elle satisfait
\[
4^\circ)\quad u_E \text{ est unimodulaire}
\]

\uline{Scholie} En \unsure{étudiant|étendant} ces objets, on étudie simultanément tous les $n$-polygônes réguliers \inserted{épinglés} pour $n\in\unsure{\mathbb N'|\mathbb N|\mathbb N^{*}},\ n\geq 3$ -- ils correspondent aux cas où $u_E$ est d'ordre fini, et $n$ est l'ordre de $u_E$ -- \unsure{en|et} même temps qu'un nouveau type d'objets : les $\infty$-polygônes réguliers épinglés.

\uline{Remarque} Les considérations générales concernant les ``polygônes réguliers'' s'étendent aussi aux $\infty$-polygônes réguliers sans chgt. ! Notamment les prop. 1, et corollaire, la prop 2 (qui \deleted{\unsure{est}}\inserted{s'étend en l'énoncé évident} concernant la \deleted{du} question de fidélité $D_\infty\hookrightarrow\mathrm{Aff}(E)$ exprimant \unsure{qu'on|que} \deleted{\unsure{conditions que|conditions}} \unsure{l'on.|bien.|finies} bien un $\infty$-poly. rég-
% slanted
\margin{! Mais il faut commencer par donner la définition des polyèdres réguliers en général (finis ou pas), en \deleted{\unsure{répét.}}\unsure{reprenant} la définition d'abord donnée dans le cas fini \ldots}
% archive sheet number: 31
\scanpage{33}
la prop. 3, \deleted{\unsure{et 2}} montrant le lien exact entre $\infty$-polygônes réguliers et repr. affines du groupe diédral infini, enfin la prop 4 (où dans b), \inserted{et c),} il ne faut plus parler de questions de divisibilité, puisque $n$ est ici remplacé par $+\infty$\ldots)

\uline{NB} Ici encore, pour une classe d'isom. de réal. géom. \inserted{régulières} \deleted{\unsure{de|des}} polygône combinatoire, \unsure{déterminé|donné}, \inserted{\unsure{et les assoc.,}} \deleted{épinglée \unsure{(biff.)} il y a \unsure{bien}} un ``module'' $\alpha\in k$, -- maintenant $\alpha\in k$ n'est plus soumis à aucune condition, cyclotomique ou autre --) il y a lieu d'écrire ``\uline{la}'' \unsure{espa.|repr.} affine $E$ correspondant \uline{au} représentant épinglé (aut. unique) Plus \unsure{gén$^{\rm t}$|général} et plus joliment ou intrinsèque, si $\Pi$ est un poly combinatoire fini ou infini, \unsure{notant|notons} que la cat. de ses réalisations géom. régulières est \uline{rigide} (ses objets, n'ont d'auto. autres que l'identité) il y a lieu pour \unsure{la} classe d'iso (déterminée tjrs par le $=\alpha$) d'écrire $E\inserted{=E(\Pi)}$ -- \deleted{étant entendu} que ce sera une
% slanted
\margin{C'est un des intérêts techniques d'éliminer \deleted{\unsure{au départ}} au début un \unsure{choix} de la théorie\ldots}
% archive sheet number: 32
\scanpage{34}
écriture intrinsèque, de façon que l'opération de $\Gamma\unsure{\Rightarrow|=|\simeq}\deleted{\operatorname{Aut}}(\Pi)$ sur $E(\Pi)$ sera obtenue automatiquement. (Il restera alors simplement à décrire $S\longrightarrow E(\Pi)$ compatible avec l'op. de $\Gamma$\ldots).

$\overline{\underline{\mathrm{VI}}}$ \uline{Opérations} $\psi_r$ \uline{d'étoilage d'un polygône.}

Soit $(S,\omega)$ un polygône combinatoire, et soit $r\in\mathbb Z$, considérons
\[
\psi_r(\omega)=\{u^r\mid u\in\omega\}\tag*{(23)}
\]
\deleted{\unsure{S'il y a|A' que}}

\uline{Prop 6} a) Pour que \deleted{$\unsure{\psi_r(\omega)}$} $(S,\psi_r(\omega))$ soit un polygône combinatoire (i.e. pour que $u^r$ opère transitivement sur $S$) il f. et s. que \deleted{$S$} $\operatorname{Card}(S)=n$ soit fini, et que $(r,n)=1$. \deleted{Si $\operatorname{card}(S)=n$,}

\inserted{b) On a, pour $r,r'\in\mathbb Z$}
\[
\inserted{\psi_r(\psi_{r'}(\omega))=\psi_{rr'}(\omega)}
\]

c) Si \deleted{$\operatorname{card}S=$} $n$ est fini, alors \deleted{\unsure{l'on a}} \deleted{\unsure{pour les} $\psi_r(\omega)$} et \unsure{soit|pour} $r,r'\in\mathbb Z$, \deleted{\unsure{on a 1)}} alors on a
\[
\psi_r(\omega)=\psi_{r'}(\omega)\quad\text{ssi}\quad r'\equiv\pm r\ (n)
\]
\inserted{donc $\Psi_n$} \deleted{\unsure{donne}} \deleted{$r\mapsto\psi_r(\omega)$ \unsure{induit une application injective} $\unsure{\to}$} \deleted{$\mathbb Z_n^*/\pm\longrightarrow\operatorname{Pol}(S)$} \inserted{opère librement sur $\operatorname{Pol}(S)$}

i.e. \deleted{\unsure{que}} $(\mathbb Z/n\mathbb Z)^*$ \uline{opère} sur l'ens $\operatorname{Pol}(S)$ des structures polygonales sur $S$. En fait, c'est
% margin note written sideways
\margin{(24) $(\mathbb Z/n\mathbb Z)^*/\{\pm1\}\mathrel{\underset{\mathrm{déf}}{=}:}\Psi_n$ qui opère sur $\operatorname{Pol}(S)$, car $-1$ opère de façon triviale}
% archive sheet number: 33
\scanpage{35}
\deleted{\unsure{la même}} Soit $u\in\omega$. La condition que $(S,\omega)$ soit un poly comb. est que $\mathbb Z$ opère (via $u$) transitivement sur $S$, \inserted{et $n\geq3$,} ou encore que $S$ est un torseur sous $\mathbb Z_n$ (où $\mathbb Z_n=\mathbb Z/n\mathbb Z$ si $n$ fini, $\mathbb Z_\infty=\mathbb Z$). \deleted{\unsure{Dans la question|Dans l'opération de}} \inserted{Or l'opération de} $\mathbb Z$ sur $S$ déduite de $u^r\in\psi_r(\omega)$ n'est autre que celle qui se factorise par $\mathbb Z_n$ et se déduit de l'op. précédente via l'hom
\diagram{\begin{tikzcd}[ampersand replacement=\&]
\mathbb Z_n \arrow[r,"r\,\mathrm{id}_{\mathbb Z_n}"] \& \mathbb Z_n
\end{tikzcd}}
Donc elle fait de $S$ un $\mathbb Z_n$-torseur ssi $r\,\mathrm{id}_{\mathbb Z_n}$ est un iso, ce qui signifie que $n$ fini et $(r,n)=1$.

(b) est trivial \deleted{\unsure{Si} $(\unsure{r})$} \deleted{On a $\psi_r(\omega)=\psi_{r'}(\omega)$}

\deleted{\unsure{ssi}} (c) $\psi_r(\omega)=\psi_{r'}(\omega)$ signifie
\[
u^{r'}=u^r\quad\text{ou}\quad u^{r'}=(u^{-1})^r\ (=u^{-r})\quad\text{donc}
\]
\[
r'\equiv r\ (n)\quad\text{ou}\quad r'\equiv-r\ (n).
\]

\uline{Remarque} Une structure de $n$-pol. \uline{orienté} sur $S$ équivaut à la donnée d'une permutation circulaire $u$ sur $S$,
% archive sheet number: 34
\scanpage{36}
ou encore : la donnée d'une structure de $\mathbb Z_n$-torseur sur $S$ ($n=\operatorname{card}S$).

On trouve l'opération $\psi_r$ ($r\in\mathbb Z_n^{\unsure{\times|*}}$) \deleted{sur}
\[
\psi_r(u)=u^r\tag*{(25)}
\]
\deleted{sur} l'ens des structures pol. orientées sur $S$, donc on a une opération de $\mathbb Z_n^\times$ sur l'ens $\operatorname{Polor}(S)$ des structures de polygône orienté sur $S$. Ici cette opération est encore libre.

Si \inserted{$P=(E,S,\omega)$} est un $n$-polygône régulier, on définit \deleted{\unsure{un}}
\[
\psi_r(E)=(E,S,\psi_r(\omega))\qquad(\text{si }(r,n)=1)\tag*{(26)}
\]
qui a $=$ espace affine \inserted{$E$ sous-jacent} et $=$ ens de sommets \deleted{\unsure{en}} \unsure{tant} que partie de $E$, mais avec une autre structure combinatoire dessus. Ici \deleted{$(\mathbb Z/n\mathbb Z)^{\unsure{*}}/\pm1$}$\Psi_n$ opère librement sur l'ens
\[
\operatorname{Polreg}_n(E)
\]
des $n$-polygônes réguliers plongés \unsure{ds|de} $E$. De $=$ $\mathbb Z_n^\times$ opérera \inserted{librem$^{\rm t}$} sur l'ens des $n$-polyg. rég. orientés. Item pour pol. rég. \uline{épinglés}
% archive sheet number: 35
\scanpage{37}
(\unsure{interprétés|plutôt} comme $n$-poly. rég. orientés munis d'un sommet). \deleted{Enfin, si $\Pi$ est un $n$-polygône combinatoire,} \inserted{la condition d'unimodularité est invariante.}

\uline{NB} Si on a une \deleted{repr.} réal$^{\rm on}$ géom. régulière d'un $n$-polygône comb. $\Pi=(S,\omega)$ ($n$ fini), alors $\psi_r$ sera la \unsure{représ.|repr.} de $\psi_r(\Pi)=(S,\psi_r(\omega))$ définie par la $=$ application $S\longrightarrow E$ -- la condition de régularité sera en effet invariante p.r. au passage de $\omega$ à $\psi_r(\omega)$, \deleted{\unsure{i.e. de}} étant donné que \deleted{le groupe des} l'on a évidemment
\[
\operatorname{Aut}(\Pi)=\operatorname{Aut}(\psi_r(\Pi))\tag*{(27)}
\]
en tant que ss-groupes de $\unsure{\mathfrak S_S|\mathcal S_S|G_S}$, i.e. le normalisateur dans $\mathfrak S_S$ de \inserted{$\omega$ et de} \deleted{\unsure{repr.|$\Psi_r$}} $\psi_r(\omega)$ est évidemment le $=$ (ce qui normalise $\omega$ normalise $\inserted{\omega'=}\psi_r(\omega)$, et inversemt car on a aussi $\omega=\psi_{r'}(\omega')$, si $r'$ est un inverse de $r$ dans $\Psi_n$).

% vertical rule beside the question
\uline{Question} Peut-on avoir \uline{isomorphie} entre un $n$-polygône \inserted{régulier} $P$ et $\psi_r(P)$, ? si $r\in\Psi_n$, $r\neq1$ dans $\Psi_n$ (i.e. si $r\in\mathbb Z_n^{\unsure{\times|*}}$, $r\neq\pm1$ dans $\mathbb Z_n^\times$) ? On trouvera
% archive sheet number: 36
\scanpage{38}
en termes de l'invariant $\alpha$ que \deleted{\unsure{cela}} l'isomorphie entre $P$ et $\psi_r(P)$ équivaut à une condition de la forme $\psi_r(\alpha)=\alpha$ [où $\psi_r\in\mathbb Z[\unsure{U|t}]$ est un pol. conv., \deleted{\unsure{\}}} ne dépendant que de $r$, pas de $n$\ldots]

On \deleted{\unsure{trouve|trouvera}} que \deleted{\unsure{l'on a}} si $(n,p)=1$ on aura $\psi_r(\alpha)\neq\alpha$ si $r\neq1$ dans $\Psi_n$, tandis que si \deleted{$\unsure{(n,p)\neq1|\text{complexe}}$} $p\mid n$ (donc ou bien $n=p$, $p$ impair, ou bien $n=2p$) alors \deleted{\unsure{tout|tous}} les $n$-pol. réguliers sur $k$ sont isom. \deleted{Mais} Donc dans ces cas \deleted{\unsure{on a, si}} \inserted{(où on a} $(\mathbb Z/p\mathbb Z)^*=\mathbb Z/(p-1)\mathbb Z$ de card$^{\rm al}$ \inserted{$p-1$} $>2$ \inserted{i.e. $\unsure{\Gamma_n|\Psi_n|r_n}\neq1$} sauf si $p=\unsure{3|2},3$ ( \deleted{le cas $p=2$ étant par ident$^{\rm é}$} \inserted{correspondant aux cas du triangle, du} carré et de l'hexagone) la réponse est négative, sauf dans les seuls \deleted{\unsure{cas}} trois cas envisagés (triangle \inserted{et hexagone} en car. 3, carré en car. 2).

On peut retenir quand $=$ que $\psi_r$ a une nette tendance à changer la classe d'isomorphie d'un $n$-polygône régulier.

On va \inserted{voir si on peut} définir les $\psi_r$ ($r\in\mathbb Z$) \uline{sans restriction sur $r$} pour les repr.
% archive sheet number: 37
\scanpage{39}
géom. du $\infty$-polygône combinatoire épinglé $\Pi_0(\infty)$. Pour ceci, on note que \deleted{la donnée de la} \deleted{dans le cas de $\Pi_0(n)$, et $r$ premier à $n$, on aura la repr. de} (pour $n$ quelconque, fini ou infini) la donnée de la repr.
\[
D_n\mathrel{\makebox[0pt][l]{\makebox[1.35em]{$\rho$}}\unsure{\rule{1.35em}{0.4pt}|\longrightarrow}}\mathrm{Aff}(E)
\]
équivaut aussi à la donnée de
\[
\left\{\begin{aligned}
&\sigma_E,u_E\in\mathrm{Aff}(E)\\
&\text{satisfaisant}\\
&\sigma_E^2= \mathrm{id}_E,\ \sigma_Eu_E\sigma_E^{-1}=u_E
\qquad\text{et si $n$ fini}\\[-2pt]
&\hspace{19em}u_E^n=\mathrm{id}_E
\end{aligned}\right.\tag*{(28)}
\]
tandis que la donnée suppl. $s_0\in E$ et les conditions $(\unsure{2^\circ|3^\circ})$ et $(4^\circ)$ \deleted{\unsure{(ou $3^\circ$, $4^\circ$)}} de la prop 5 \deleted{ne dépendent pas} \deleted{\unsure{survivent}} \inserted{\unsure{sont stables|détaillées} par} le chgt de $u_E$ en $u_E^r$ -- qui a un sens quel que soit $r\in\mathbb Z$, et (pour $n$ fini) continue à vérifier la condition $u_E^n=\mathrm{id}$ (après remplacement de $u_E$ par $u_E^r$). \deleted{\unsure{(Ici il}} Quant à la condition $2^\circ)$ de non align$^{\rm mt}$ de $s_0$, $u^rs_0$ et $u^{2r}s_0$ (ou de tous les $u^{ir}s_0$, $i\in\mathbb Z$, ce qui revient au $=$)
% archive sheet number: 38
\scanpage{40}
\deleted{elle} pour $r$ fixé, \deleted{elle} \inserted{le fait si elle \deleted{\unsure{est}} \unsure{est|n'est} \inserted{\unsure{ou non|bien}} satisfaite par $u_E^r$} dépend des \deleted{le} choix de $u_E$ et $s_0$ ($\sigma_E$ est ici indiffér., à vrai dire) -- donc du ``module'' $\alpha$.

\deleted{\unsure{On s'intéressera à|On a,}} P. ex. si \inserted{$n=\infty$ mais si} $\alpha$ est tel que la repr. de $D_\infty$ se factorise par $D_r$, donc $u_r=\mathrm{id}$, cette condition ne sera pas satisfaite.

Néanmoins on verra, \inserted{comme on a dit,} que $\psi_r(\alpha)$ peut être défini pour \uline{tt} $\alpha$ -- il se pose alors la question de la signification géom de \deleted{$\psi_r(\alpha)$} la repr. de $D_{\deleted{\unsure{\infty|n}}}$ associé à $\psi_r(\alpha)$, dans le cas ``exceptionnel'' [à caractériser en termes de $\alpha,r$ ! \inserted{(c'est subordonné au Pb général, pour $\alpha$ fixé, de savoir quand trois sommets $s_i,s_j,s_k$ sont alignés\ldots)}] où les $s_{ir}=u^{ir}s_0$ sont tous alignés \ldots.
% slanted
\margin{Question à élucider \ldots !}

$\overline{\underline{\mathrm{VII}}}$ \deleted{Polygônes} \inserted{\uline{Systèmes} \deleted{\unsure{de}}} \uline{réguliers} \deleted{\unsure{finis}} \inserted{\uline{de $n$ sommets}}

% vertical rule beside the parenthesis
(i.e. structures polygonales modulo étoilage \ldots)

Soit $S$ un ensemble \inserted{fini, de cardinal $n$}. On appelle \uline{structure quasi-polygonale} \inserted{(combinatoire)} sur $S$ la donnée d'une structure polygonale
% archive sheet number: 39
\scanpage{41}
sur $S$ modulo une des opérations $\psi_r$ ($r\in\mathbb Z_n^{*}$, ou $r\in\Psi_n$) ; \deleted{\unsure{Il revient donc au $=$|Je dis}} \inserted{en d'autres} termes, c'est une orbite $\pi$ de $\Psi_n$ opérant sur $\operatorname{Pol}(S)$. Regardant le ss-groupe \inserted{$Z$ de $\unsure{\mathfrak S_S|\mathcal S_S|G_S}$} engendré par un $\omega\in\pi$ (lequel ne dépend que de $\pi$, \deleted{\unsure{on est en|étant}} \inserted{étant} invariant en remplaçant $\omega$ par $\psi_r(\omega)$), on trouve un ss-groupe cyclique d'ordre $n$, et $S$ est un torseur sous ce groupe. Ainsi, \inserted{la donnée d'}une structure quasi-polygonale \deleted{sur} $S$ équivaut à la donnée d'une structure de torseur sur $S$ à groupe de \deleted{\unsure{gr|des}} translations \inserted{$Z$} cyclique (d'ordre $n$ nécessairement). Se donner un $\omega\in\pi$ revient au $=$ que de se donner un gén. ``au signe près'' dans $Z$. Se donner une structure polygonale \uline{orientée} compatible avec la structure quasi-polygonale revient à la donnée d'un gén. de $\unsure{Z|\mathbb Z}$.

La catégorie des \deleted{quasi} quasi-polygônes combinatoires d'ordre $n$ ($n\in\mathbb N$, $n\geq3$) est équivalente à celle des fibrés en droites affines sur $\operatorname{Spec}(\mathbb Z/n\mathbb Z)$, à un tel
% archive sheet number: 40
\scanpage{42}
fibré correspondant l'ens de ses sections, qui est un torseur sur le groupe \inserted{$Z$} des sections du fibré vectoriel correspondant, lequel est un $\mathbb Z/n\mathbb Z$-module libre de rang 1, i.e. un groupe cyclique d'ordre $n$.

Si $S$ est un quasi-pol. combinatoire, $Z$ son groupe des rotations, le groupe de ses automorphismes \inserted{(noté $\mathrm{Aff}(S)$)} est canoniquement muni d'une structure d'extension
\diagram{\begin{tikzcd}[ampersand replacement=\&,column sep=small,row sep=small]
0\arrow[r] \& Z\arrow[r] \& \mathrm{Aff}(S)\arrow[r] \& (\mathbb Z/n\mathbb Z)^{*}\arrow[r]\arrow[d,phantom,"\rotatebox{90}{$\simeq$}" description] \& 0\\
\& \& \& \operatorname{Aut}(Z) \&
\end{tikzcd}}
et le choix d'un $s\in S$ détermine un splitting de cette extension, permettant d'identifier $\mathrm{Aff}(S)$ au produit $\frac12$ direct $(\mathbb Z/n\mathbb Z)^*\mathbin{.}Z$. L'image inverse du ss-grpe $\pm1$ de $(\mathbb Z/n\mathbb Z)^*$ dans $\mathrm{Aff}(S)$ est un ss-grpe $\mathrm{Aff}^{\pm}(S)$ de $\mathrm{Aff}(S)$ d'ordre $2n$. C'est aussi le
% margin note written sideways
\margin{ss-groupe $\operatorname{Aut}(S,\omega)$ de $\unsure{\mathfrak S_S|\mathcal S_S|G_S}$, où $\omega$ est n'importe laquelle des str. pol. sur $S$ compatibles avec sa str. quasi-pol.}

Soit $E$ un plan affine sur corps $k$.

Un \uline{système régulier de sommets} dans $E$ est \deleted{\unsure{la co}} un couple $(S,Z)$, où $S\subset E$ est une partie finie,
% archive sheet number: 41
\scanpage{43}
$Z$ est une structure quasi-régulière sur $S$ (identifiée p.ex. à un ss-groupe cyclique de $\unsure{\mathfrak S_S|\mathcal S_S|G_S}$ d'ordre $n=\operatorname{card}(S)$), tels que pour une (donc tte) structure polygonale $\omega$ sur $S$ compatible avec $Z$ (i.e. telle que $\omega\subset Z$), $(E,S,\omega)$ soit un polygone régulier. Il revient au $=$ de dire que tout $u\in\mathrm{Aff}^{\pm}(S,Z)$ se prolonge de façon unique en un autom. affine de $E$. [NB On ne dit rien sur les autres automorphismes de $\mathrm{Aff}(S,Z)$ \deleted{\unsure{ca}}
% slanted
\margin{Question à élucider}
(il y en a bien d'autres, sauf dans le cas où $\unsure{\varphi(n)|\phi(n)}=2$ \inserted{nb d'Euler} i.e. $n=\deleted{2,}3,4$ ou $6$)] On peut encore dire que \inserted{la donnée d'}un \deleted{\unsure{gr|un}} syst. rég. de $n$ sommets dans $E$ équivaut à la donnée d'une orbite de $\Psi_n\inserted{=(\mathbb Z/n\mathbb Z)^{\unsure{*|\times}}/\{\pm1\}}$ opérant sur l'ens de tous les $n$-polygônes réguliers dans $E$.
% archive sheet number: 42
\scanpage{44}
\uline{Question} Si $(S,Z)$ est un système $n$-régulier de sommets dans $E$, $Z$ est-il déterminé de façon unique par la donnée de la seule partie $S$ de $E$ ? On trouvera ceci : si $(S,\omega)$ est un $n$-polygône régulier dans $E$, le stabilisateur \inserted{$G$} de $S$ dans $\mathrm{Aff}(E)$ (je dis bien de \inserted{la \uline{partie}} $S$ de $E$, non du couple $(S,\omega)$) est égal au groupe $\operatorname{Aut}(E,S,\omega)$, sauf dans les cas $\alpha=2$, \deleted{où on trouve un groupe isom} \inserted{ou $\alpha=-2$ ; dans ces cas, identifiant} $G$ à un sous-groupe de $\unsure{\mathfrak S_S|\mathcal S_S|G_S}$, on trouve $\mathrm{Aff}(S,Z)$.
% slanted
\margin{\uline{à vérifier} !}

Donc la question devient purement combinatoire, savoir : une structure quasi polygonale comb. \inserted{$Z$} sur $S$ est-elle déterminée par la seule connaissance de $\mathrm{Aff}^{\pm}(S)\subset\mathfrak S_S$, ou par celle de $\mathrm{Aff}(S)\subset\mathfrak S_S$ ? Dans le premier cas, $Z$ est déterminé comme le plus grand ss-groupe cyclique \uline{distingué} $\neq G=\mathrm{Aff}^{\pm}(S)$. [car si $\sigma\in G\smallsetminus Z$, on voit facilem$^{\rm t}$ que
% archive sheet number: 43
\scanpage{45}
le groupe distingué \deleted{qu'il} engendré par $\sigma$ \inserted{contient les rotations paires de $Z$} \deleted{donc est \unsure{indice}} \inserted{égal à $G$} si $n$ impair, et sinon \unsure{un|il} \deleted{\unsure{groupe}} \deleted{\unsure{se}} est isomorphe à $\{\pm1\}\mathbin{.}\mathbb Z/n'\mathbb Z$ (produit $\frac12$ direct) où $n'=n/2$ -- donc il est non cyclique, donc $\sigma$ n'est pas contenu dans un ss-groupe invariant cyclique $\neq G$, donc tous ces groupes sont contenus dans $Z$, qui est \inserted{donc} le plus grand de ces groupes.] Dans le cas de $\mathrm{Aff}(S)$\ldots.
% slanted
\margin{faire dém\ldots}

Donc sauf erreur, la structure quasi-polygonale combinatoire de $S\inserted{\subset E}$ est \inserted{uniquemt} déterminée, pour son système régulier de $n$ sommets dans $E$.

Principe de classification géométrique, pour $n$ fixé : Un système régulier de $n$ sommets (\deleted{\unsure{un|en}}\inserted{un} $\mathrm{SR}\deleted{\unsure{S}}n\mathrm S$) détermine

a) Un sous groupe cyclique d'ordre $n$
\[
Z\subset\mathrm{Aff}(E)
\]

b) Une orbite $\displaystyle S\in {}_Z\!\backslash E$

avec les conditions suivantes

1) $S$ non contenu dans une droite (cela implique que $Z$ opère \uline{fidèlemt} sur $S$, qui est donc un $Z$-torseur, $Z$ étant cyclique)

2) $\exists\,\sigma\in\operatorname{Norm}(Z)$ \deleted{$\unsure{\mathrm{Aff}|{}\neq\mathrm{id}_E}$} tel que $\sigma$ \deleted{\unsure{laisse fixe un pt de $S$|laisse invariant un pt de $S$}} \inserted{opère sur $Z$ par $z\mapsto z^{-1}$, et que $\sigma$ \unsure{invarie|inverse|invariant}} un pt de $S$ --
% archive sheet number: 44
\scanpage{46}
\deleted{Réciproquement, si ces conditions sont satisfaites, choisissant un gén $u$ de $Z$, on voit que} -- de sorte que $\{1,\sigma\}\mathbin{.}Z$ (produit $\frac12$ direct) opère sur $S$,\ldots \deleted{qui devient un polygône de}

On a évidemment

\uline{Prop 7} La donnée d'un \deleted{$n$-polygône} $\mathrm{SR}n\mathrm S$ dans $E$ équivaut à la donnée de a) b) ci-dessus (ss-grou \inserted{$Z$} cyclique d'ordre $n$ de $\mathrm{Aff}(E)$, et d'une orbite $S$ \deleted{de} $Z$ dans $E$). \deleted{faisant de \unsure{($E$ un type|$E$ un syst.)|les conditions}} \inserted{tels que les cond$^{\rm s}$} $1^\circ)$ et $2^\circ)$ ci-dessus soient vérifiées ($S$ non sur une droite, et condition de normalisation\ldots) -- et éventuellement la condition unimodulaire :

3) $Z$ formé d'endom unimodulaires.

\uline{Remarque} Suivant les cas ($p=\mathrm{car}\, k$, \unsure{etc|et|$n$)}, le groupe des autom d'un \deleted{syst} $\mathrm{SR}n\mathrm S$ est isomorphe à $\mathrm{Aff}^{\pm}(S)\simeq D_n$, ou à $\mathrm{Aff}(S)$. On peut se poser la question, dans l'un ou l'autre cas, de \deleted{savoir}
% archive sheet number: 45
